% généré par scripts/04_generated.py
\begin{longtable}{@{}r p{0.70\linewidth} r@{}}
\caption{Collèges consulaires 340–403 selon la tradition des fastes et nombre de textes distincts (ensemble B) classés sous chaque année. Table non collationnée sur la CLRE.}\label{tab:fastes}\\
\toprule Année & Collège consulaire & Textes\\\midrule\endfirsthead
\toprule Année & Collège consulaire & Textes\\\midrule\endhead
\bottomrule\endfoot
340 & Acindynus et Proculus & 3 \\
341 & Marcellinus et Probinus & 1 \\
342 & Constantius Aug. III et Constans Aug. II & 2 \\
343 & Placidus et Romulus & 3 \\
344 & Leontius et Sallustius & 1 \\
345 & Amantius et Albinus & 3 \\
346 & Constantius Aug. IV et Constans Aug. III (à Rome : p.c. Amantii et Albini) & 4 \\
347 & Rufinus et Eusebius & 3 \\
348 & Philippus et Salia & 8 \\
349 & Limenius et Catullinus & 4 \\
350 & Sergius et Nigrinianus & 2 \\
351 & Occident (Magnence) : Magnentius Aug. et Gaiso ; Orient : p.c. Sergii et Nigriniani & 1 \\
352 & Occident (Magnence) : Decentius Caes. et Paulus ; Constance : Constantius Aug. V et Constantius Caes. & 4 \\
353 & Occident (Magnence) : Magnentius Aug. II et Decentius Caes. II ; Constance : Constantius Aug. VI et Constantius Caes. II & 1 \\
354 & Constantius Aug. VII et Constantius Caes. III & 0 \\
355 & Arbitio et Lollianus & 4 \\
356 & Constantius Aug. VIII et Iulianus Caes. & 2 \\
357 & Constantius Aug. IX et Iulianus Caes. II & 2 \\
358 & Datianus et Cerealis & 5 \\
359 & Eusebius et Hypatius & 2 \\
360 & Constantius Aug. X et Iulianus Caes. III & 3 \\
361 & Taurus et Florentius & 3 \\
362 & Mamertinus et Nevitta & 11 \\
363 & Iulianus Aug. IV et Sallustius & 10 \\
364 & Iovianus Aug. et Varronianus & 9 \\
365 & Valentinianus Aug. et Valens Aug. & 6 \\
366 & Gratianus nob. puer et Dagalaifus & 7 \\
367 & Lupicinus et Iovinus & 18 \\
368 & Valentinianus Aug. II et Valens Aug. II & 7 \\
369 & Valentinianus nob. puer et Victor & 4 \\
370 & Valentinianus Aug. III et Valens Aug. III & 8 \\
371 & Gratianus Aug. II et Probus & 12 \\
372 & Modestus et Arintheus & 11 \\
373 & Valentinianus Aug. IV et Valens Aug. IV & 6 \\
374 & Gratianus Aug. III et Equitius & 7 \\
375 & Pas de consuls ordinaires : p.c. Gratiani III et Equitii & 10 \\
376 & Valens Aug. V et Valentinianus iun. Aug. & 6 \\
377 & Gratianus Aug. IV et Merobaudes & 11 \\
378 & Valens Aug. VI et Valentinianus iun. Aug. II & 3 \\
379 & Ausonius et Olybrius & 3 \\
380 & Gratianus Aug. V et Theodosius Aug. & 6 \\
381 & Syagrius et Eucherius & 10 \\
382 & Antonius et Syagrius & 13 \\
383 & Merobaudes II et Saturninus & 20 \\
384 & Richomeres et Clearchus & 15 \\
385 & Arcadius Aug. et Bauto & 6 \\
386 & Honorius nob. puer et Euodius & 8 \\
387 & Valentinianus Aug. III et Eutropius & 4 \\
388 & Theodosius Aug. II et Cynegius (Occident : Magnus Maximus Aug. II) ; Merobaudes III attesté par l'épigraphie (n° 298) & 10 \\
389 & Timasius et Promotus & 13 \\
390 & Valentinianus Aug. IV et Neoterius & 10 \\
391 & Tatianus et Symmachus & 13 \\
392 & Arcadius Aug. II et Rufinus & 8 \\
393 & Theodosius Aug. III et Abundantius ; à Rome : Eugenius Aug. & 11 \\
394 & Arcadius Aug. III et Honorius Aug. II ; territoire d'Eugène : Nicomachus Flavianus & 6 \\
395 & Olybrius et Probinus & 17 \\
396 & Arcadius Aug. IV et Honorius Aug. III & 10 \\
397 & Caesarius et Atticus & 22 \\
398 & Honorius Aug. IV et Eutychianus & 12 \\
399 & Mallius Theodorus et Eutropius (Occident : Theodorus seul) & 9 \\
400 & Stilicho et Aurelianus & 8 \\
401 & Vincentius et Fravitta & 13 \\
402 & Arcadius Aug. V et Honorius Aug. V & 9 \\
403 & Theodosius iun. Aug. et Rumoridus & 10 \\
\end{longtable}
